% Example a=1, lambda=2; positive signed projection along C->A.
\begin{center}\begin{tikzpicture}[font=\small]
\draw[black!30,rounded corners=2pt] (0,0) rectangle (14.4,8.2);
\node[anchor=west,font=\small\bfseries] at (.5,7.72)
  {速算法先看勾股：作垂足、拆距离，再看哪些项会抵消};
\node[anchor=west] at (.8,6.8) {示例 $\lambda=2$：$CA=4CB$；射线 $CA,CB$ 同向};
\begin{scope}[shift={(4.5,3.7)},scale=2.5]
\coordinate (A) at (-1.0000000000,0.0000000000);
\coordinate (B) at (1.0000000000,0.0000000000);
\coordinate (C) at (1.6666666667,0.0000000000);
\coordinate (P) at (0.6000000000,0.8000000000);
\coordinate (H) at (0.6000000000,0.0000000000);
\draw[black!50] (-1.25,0)--(2.1,0);
\draw[blue,very thick] (A)--(P);
\draw[teal!80!black,very thick] (B)--(P);
\draw[orange!85!black,thick] (C)--(P) node[midway,above right=3pt] {$PC$};
\draw[black!55,dashed] (P)--(H) node[midway,left=5pt] {$PH$};
\draw[black!60] (.6,.12)--(.72,.12)--(.72,0);
\foreach \q in {A,B,C,P,H}{\fill (\q) circle (1.3pt);}
\node[below=6pt] at (A) {$A$};
\node[below right=6pt] at (B) {$B$};
\node[below right=6pt] at (C) {$C$};
\node[below left=6pt] at (H) {$H$};
\node[above=5pt] at (P) {$P$};
\draw[orange!85!black,->,thick] (1.6666666667,-.55)--(.6,-.55)
  node[midway,below=3pt] {$t=CH$};
\end{scope}
\node[anchor=west] at (.8,1.15)
  {先作 $PH\perp AB$；图中 $t=CH>0$，即 $CP$ 沿 $C\to A$ 方向的投影。};
\node[anchor=west] at (.8,.5)
  {先用 $PH^2=PC^2-t^2$，再分别在 $\triangle PAH,\triangle PBH$ 中用勾股定理。};
\end{tikzpicture}\end{center}
